\section{Two elementary sieve facts}\label{sec:sieve}

The prime construction needs two elementary facts. The first estimates
the mass left after excluding local divisibility conditions, with an
explicit bound on the moduli in the error term. The second identifies
the density removed when a least prime factor is selected. We give
proofs, following the forms in \cite[Lemmas~2.9 and~7.4]{BlockPaper}.

\subsection{A block sieve with bounded coefficients}

The advantage of the following form of the Brun--Hooley sieve is that
its remainders occur without divisor weights. Its underlying product
inequalities are those of Ford and Halberstam
\cite[Section~1, Lemma~1]{FordHalberstam2000}.

\begin{lemma}[Block sieve]\label{lem:block}
Let $\mathcal A$ be a finite set with nonnegative weights, and let
$\mathcal P$ be a set of primes at most $z$, where $z\ge2$.
For each $p\in\mathcal P$, specify a bad condition on $\mathcal A$.
Suppose that, for every squarefree product $d$ of primes in
$\mathcal P$, the weight of the objects satisfying all conditions
at $p\mid d$ is
\[
 A(d)=X'g(d)+E(d),
\]
including $d=1$, where $X'\ge0$, $g(1)=1$, and $g$ is multiplicative.
Assume that, for fixed constants $\eta_0>0$ and $C_0<\infty$,
\begin{equation}\label{eq:block-hypotheses}
 0\le g(p)\le1-\eta_0,
 \qquad
 \sum_{\substack{v<p\le v^2\\p\in\mathcal P}}g(p)\le C_0
 \quad(v>1).
\end{equation}
For every sufficiently large even integer $H$, the weight $S$ of
objects avoiding all the bad conditions satisfies
\begin{equation}\label{eq:block-sieve}
 S=X'\prod_{p\in\mathcal P}(1-g(p))
       \bigl(1+O(\exp(-H))\bigr)
   +O\left(\sum_{\substack{d\le z^{4H+2}\\d\mid\prod_{p\in\mathcal P}p}}
                  |E(d)|\right).
\end{equation}
The constants and the lower bound on $H$ depend only on
$\eta_0,C_0$. In particular, the estimate is uniform when $H$ grows.
For $H=2$ there is the upper bound
\begin{equation}\label{eq:block-sieve-upper}
 S\ll_{\eta_0,C_0}X'\prod_{p\in\mathcal P}(1-g(p))
       +\sum_{\substack{d\le z^{10}\\d\mid\prod_{p\in\mathcal P}p}}
                    |E(d)|.
\end{equation}
\end{lemma}

\begin{proof}
Write $b_p\in\{0,1\}$ for the indicator of the bad condition at $p$.
Use the blocks
\[
 \mathcal B_j=\{p\in\mathcal P:
        z^{2^{-j-1}}<p\le z^{2^{-j}}\},\qquad j\ge0,
\]
ignoring empty blocks and retaining the indicated indices.
Put $h_j=(j+1)H$. Let $I_j$ be the indicator of avoiding the
conditions in $\mathcal B_j$. Define $U_j$ to be the
inclusion--exclusion polynomial through degree $h_j$, and $T_j$ to
be the sum of the unsigned degree-$h_j+1$ inclusion--exclusion terms:
\[
 U_j=\sum_{\substack{J\subseteq\mathcal B_j\\|J|\le h_j}}
             (-1)^{|J|}\prod_{p\in J}b_p,
 \qquad
 T_j=\sum_{\substack{J\subseteq\mathcal B_j\\|J|=h_j+1}}
             \prod_{p\in J}b_p.
\]
If exactly $t$ conditions in the block hold, then, since $h_j$ is even,
\[
 U_j=\sum_{a=0}^{h_j}(-1)^a\binom ta
 =\begin{cases}1,&t=0,\\ \binom{t-1}{h_j},&t\ge1.
 \end{cases}
\]
As usual, a binomial coefficient vanishes when its lower index
exceeds its nonnegative upper index. Thus
$0\le I_j\le U_j$ and $U_j-I_j\le T_j$.
Telescoping the product, all of whose factors are nonnegative, gives
\[
 0\le\prod_jU_j-\prod_jI_j
       \le\sum_jT_j\prod_{k\ne j}U_k.
\]
Consequently the polynomials
\begin{equation}\label{eq:block-polynomials}
 U=\prod_jU_j,
 \qquad
 B=U-\sum_jT_j\prod_{k\ne j}U_k
\end{equation}
satisfy $B\le\one_{\text{no bad condition}}\le U$.

Every monomial in $U$ has degree at most $h_j$ in block $j$.
A monomial in correction $j$ has degree exactly $h_j+1$ there
and degree at most $h_k$ in every other block. Hence the supports
of the corrections are disjoint from one another and from the
support of $U$. All coefficients of both $B$ and $U$ therefore
have absolute value at most one. The prime product indexing a
monomial of $U$ has logarithm at most
\[
 \log z\sum_{j\ge0}(j+1)H2^{-j}=4H\log z.
\]
A correction adds at most $\log z$, so every monomial used is
supported on a squarefree $d\le z^{4H+2}$.

It remains to evaluate the main terms of these polynomials.
Give the indicators $b_p$ independent Bernoulli laws with means
$g(p)$, and write $\mathbb E$ for expectation under this auxiliary
law. By \eqref{eq:block-hypotheses}, each block has total mean at
most $C_0$. Its avoidance probability satisfies
\[
 a_j:=\mathbb E I_j=\prod_{p\in\mathcal B_j}(1-g(p))
       \ge \exp(-C_0/\eta_0)=:a_*>0,
\]
because $-\log(1-u)\le u/\eta_0$ for $0\le u\le1-\eta_0$.
Moreover,
\[
 t_j:=\mathbb E T_j\le\frac{C_0^{h_j+1}}{(h_j+1)!},
 \qquad a_j\le\mathbb E U_j\le a_j+t_j.
\]
The factorial bound implies
\[
 R_H:=\sum_j\frac{t_j}{a_j}\ll_{\eta_0,C_0}\exp(-H)
\]
for all sufficiently large even $H$: for sufficiently large $n$,
$C_0^n/n!\le\exp(-2n)$, and $h_j=(j+1)H$.
For $H=2$, the same sum is bounded in terms of $\eta_0,C_0$.
Independence of the blocks now yields
\[
 \prod_ja_j\le\mathbb E U\le \prod_ja_j\exp(R_H),
 \qquad
 \mathbb E(U-B)\le \prod_ja_j R_H\exp(R_H).
\]
Thus both bounding polynomials have expectation
$\prod_ja_j(1+O(\exp(-H)))$ when $H$ is sufficiently large;
for $H=2$, the upper expectation is $O(\prod_ja_j)$.

Finally evaluate $U$ and $B$ in the original weighted set.
Multiplicativity makes their main terms $X'\mathbb E U$ and
$X'\mathbb E B$. The coefficient and support bounds show that
each error is at most the sum of $|E(d)|$ in
\eqref{eq:block-sieve}. Squeezing $S$ between the two weighted
polynomial sums proves both assertions.
\end{proof}

\subsection{Rough-number densities}

Recall the density $D_\gamma(s)$ from
\eqref{eq:rough-density-definition}. Its $j$th term describes a
product of $j$ primes whose logarithms, in units of $\log x$, sum
to $s$ and are at least $\gamma$; the factor $1/j!$ normalizes the
ordered prime lists. These are Buchstab's rough-number densities,
expressed in logarithmic coordinates. The identity below is the
corresponding least-prime-factor decomposition \cite{Buchstab1937}.
The normalization and regularity were established in
Section~\ref{sec:type-ii}; we now prove the two estimates needed to
subtract composite mass.

For the product $V(y)$ in \eqref{eq:sieve-basic-parameters},
Mertens' theorem gives
\begin{equation}\label{eq:mertens-product}
 V(y)\sim\frac{\exp(-\gamma_E)}{\log y},
\end{equation}
where $\gamma_E$ is Euler's constant.

\begin{lemma}[Rough-number density]\label{lem:density}
For $0<b<1/2$,
\begin{equation}\label{eq:density-identity}
 \int_b^{1/2}D_t(1-t)\,\frac{dt}{t}=D_b(1)-1,
\end{equation}
where the integrand at $1/2$ is interpreted by its left limit.
There are absolute constants $b_0,c_d,C_d>0$ such that, for
$0<b<b_0$,
\begin{equation}\label{eq:density-upper}
 D_b(1)\le\frac{\exp(-\gamma_E)}{b}
                 \bigl(1+C_d\exp(-c_d/b)\bigr).
\end{equation}
One may take \(c_d=1/20\).
\end{lemma}

\begin{proof}
To prove \eqref{eq:density-identity}, write a term of $D_b(1)$ with
$j\ge2$ as an integral on
\[
 t_1+\cdots+t_j=1,\qquad t_i\ge b,
 \qquad
 \frac1{j!}\frac{dt_1\cdots dt_{j-1}}{t_1\cdots t_j}.
\]
This measure is invariant under permutations of the $j$ coordinates:
interchanging a dependent and an independent coordinate has absolute
Jacobian one. Except on a set of measure zero, the minimum coordinate
is unique. Choose it in $j$ ways and call its value $t$.
The remaining $j-1$ coordinates sum to $1-t$ and are at least $t$;
their measure is the $(j-1)$-factor term of $D_t(1-t)$, multiplied
by $dt/t$, since $j/j!=1/(j-1)!$. Necessarily $t\le1/2$.
Summing over $j\ge2$ proves the identity; the omitted one-prime
term of $D_b(1)$ equals $1$.

For \eqref{eq:density-upper}, we compare $D_b(1)$ with ordinary
rough integers and then apply Lemma~\ref{lem:block}. Fix $b>0$,
put $L=\log x$, and let
\[
 R_b(x)=\#\{n\in\mathbb{N}:x<n\le2x,\ P^-(n)>x^b\}.
\]
We first claim that
\begin{equation}\label{eq:density-rough-lower}
 D_b(1)\le\liminf_{x\to\infty}\frac{L}{x}R_b(x).
\end{equation}
For $j\ge2$ and fixed $\eta>0$, count ordered prime lists
$(p_1,\ldots,p_j)$ with weight $1/j!$, subject to
\[
 p_i>x^b\quad(i<j),\qquad
 Q:=p_1\cdots p_{j-1}<x^{1-b-\eta},\qquad
 x/Q<p_j\le2x/Q.
\]
Every product counted is in $R_b(x)$, because $p_j>x^{b+\eta}$.
An integer with prime multiplicities $e_1,e_2,\ldots$ receives
total weight at most $1/\prod_i e_i!\le1$, even when prime factors
repeat. Its total number of prime factors also determines $j$.
Thus summing these counts over $j$, together with the one-prime
contribution, gives a lower bound for $R_b(x)$.

Uniformly in the indicated range of $Q$, the prime number theorem gives
\[
 \#\{p_j:x/Q<p_j\le2x/Q\}
  =(1+o(1))\frac{x}{Q\log(x/Q)}.
\]
For fixed $0<\alpha<\beta$, Mertens' theorem gives
\[
 \sum_{x^\alpha<p\le x^\beta}\frac1p
       \longrightarrow\log(\beta/\alpha).
\]
Consequently the reciprocal-prime measures on any fixed compact
interval of positive logarithmic exponents converge to $dt/t$.
Their finite products converge as well. On the truncated simplex
$t_i\ge b$, $\sum_{i<j}t_i<1-b-\eta$, the function
$(1-\sum_{i<j}t_i)^{-1}$ is bounded, and its boundary has measure
zero. The weighted count above, multiplied by $L/x$, therefore
tends to
\[
 \frac1{j!}
 \int_{\substack{t_i\ge b\ (i<j)\\
                       \sum_{i<j}t_i<1-b-\eta}}
       \frac{1}{1-\sum_{i<j}t_i}\prod_{i<j}\frac{dt_i}{t_i}.
\]
There are only finitely many relevant $j$ for fixed $b$.
The one-prime contribution tends to $1$. First let $x$ tend to
infinity with $b,\eta$ fixed, and then let $\eta$ decrease to zero.
The boundary hyperplanes are null, so these integrals increase to
the terms in $D_b(1)$. This proves \eqref{eq:density-rough-lower}.

For the upper bound on $R_b(x)$, apply Lemma~\ref{lem:block} to
the integers in $(x,2x]$, with bad conditions $p\mid n$ for
$p\le z=x^b$. Here $X'=x$, $g(d)=1/d$, and $E(d)=O(1)$,
uniformly in $d$. The hypotheses hold with absolute constants:
$g(p)\le1/2$, and the sums of $1/p$ over $(v,v^2]$ are bounded.
Choose
\[
 H=2\left\lfloor\frac{1}{40b}\right\rfloor.
\]
For sufficiently small $b$, this is an admissible even integer,
$H\gg1/b$, and $b(4H+2)\le1/5+2b<1$.
Lemma~\ref{lem:block} gives, with absolute implied constants,
\[
 R_b(x)\le xV(x^b)\bigl(1+O(\exp(-c/b))\bigr)
                +O\bigl(x^{1/5+2b}\bigr)
\]
with \(c=1/20\), since \(\exp(-H)\le\e^2\exp(-1/(20b))\).
Multiplying by $L/x$, applying \eqref{eq:mertens-product}, and
letting $x$ tend to infinity proves \eqref{eq:density-upper}
by \eqref{eq:density-rough-lower}.
\end{proof}

The term $1$ in \eqref{eq:density-identity} is the one-prime
contribution. The integral accounts for the composite contributions
by selecting their least prime factor. This is the cancellation
that will leave a positive prime mass in the construction below.
